% page 610 of the facsimile (image p-609 of the scan)
% block 152.4 x 242.0 mm ; left margin 8.0 mm
\pgc{-12.700mm}{0.000mm}{
\l{\marge[76.583mm]{24.051mm}{\ornamento{12.241mm}{ornaments/letters/letter-U.png}}}
\l{\cel{2}602}
\l{}
\l{}
\l{}
\l{}
\l{}
\l{}
\l{}
\l{ube.\cel{2}(adverbo).\cel{2}En ta loko ; en qua loko ? - FL.}
\l{}
\l{ubiqua.\cel{1}\sou{Qua}\cel{1}esas prezenta en\cel{1}\sou{omna}\cel{1}loki. EFIS.}
\l{}
\l{ucelo.\cel{2}(zool.)\cel{2}Un de la animali ek la klaso de vertebrozi,}
\l{\cel{3}ovipara, kun du pedi, la membri avana alo-forma, kovrita}
\l{\cel{3}per plumi, beko korno-harda, sen denti. - FI.}
\l{}
\l{uf !\cel{3}Interjeciono qua expresas la alejeso, la sento di li-}
\l{\cel{3}berigeso pri ulo sufriganta, peniganta, tddanta.}
\l{}
\l{\sou{ukazo.}\cel{2}I. Edikto da la caro. - II. (\sou{metaf.}) Edikto di qua}
\l{\cel{3}la autoritato esas egardenda absolute. - DEFIRS.}
\l{}
\l{\sou{-ulo.}\cel{3}Sufixo qua indikas la maskulo di la animalo. (Ta}
\l{\cel{3}sufixo uzesas nur kande lu esas nekareebla).}
\l{}
\l{\sou{ula}.\cel{2}Qua ne determinesas - nek mem esas determininda.- L.}
\l{}
\l{\sou{ulano}. (\sou{olim.})\cel{2}Soldato ek regimento de lancieri, en ula}
\l{\cel{3}landi.\cel{2}- DEFIRS.}
\l{}
\l{\sou{ulcero.}\cel{2}(\sou{patol.})\cel{2}Lezuro di tisuo qua, pro inherar la kons-}
\l{\cel{3}tituciono, tendencas adextensar su sencese. - EFIS.}
\l{}
\l{\sou{ulexo}.\cel{2}(\sou{bot.})\cel{2}Arbusto qua kreskas en la erikeyi. -L.}
\l{\cel{3}\sou{ulex}.}
\l{}
\l{\sou{ulmo}.\cel{2}(\sou{bot.)}\cel{3}Varietato di la fraxino. - L.\cel{1}\sou{ulmus}.}
\l{\cel{3}- DEFIS.}
\l{}
\l{\sou{ulno}.\cel{2}Olima mezur-unajo por la texuro. (En Paris = 1,18}
\l{\cel{3}metro). FIS.}
\l{}
\l{\sou{ultima}. Qua ne sequesos da\cel{2}nexta ; pos qua nula sama even-}
\l{\cel{3}tos. - EFIS.}
\l{}
\l{\sou{ultimato}. (\sou{diplomac.})\cel{2}Dico serio-lasta, ultima, rezolvo}
\l{\cel{3}quan onu deklaras ne-revokebla. - DEFIRS.}
\l{}
\l{\sou{ultramaro}. Farbo bele-azurea. DEFIRS.}
\l{}
\l{\sou{ultramikroskopo}. (\sou{tekn.})\cel{2}Aranjuro partikulara di lumizo,}
\l{\cel{3}adaptita a mikroskopo, qua posibligas vidar la existo}
\l{\cel{3}e la movi di ob1ekt1 d1 uui 1\cel{2}ein\cel{3}i\cel{2},}
\l{\cel{3}kam olti observebla\cel{2}ye la plugrandigo maxima di mikro-}
\l{\cel{3}skopo ordinara. - DEF.}
}
